\subsection{Analisis Kebutuhan Audit Event}

Berdasarkan hasil analisis ancaman menggunakan metode STRIDE pada bagian sebelumnya, tahap selanjutnya adalah menentukan kebutuhan audit trail yang diperlukan untuk menghasilkan bukti digital terhadap setiap ancaman yang teridentifikasi. Kebutuhan tersebut diperoleh dengan memetakan setiap ancaman ke aktivitas sistem yang perlu dicatat sebagai audit event. Hasil analisis ini menjadi dasar dalam mengidentifikasi kesenjangan mekanisme audit pada DecMed serta menjadi masukan bagi perancangan sistem audit trail pada bab selanjutnya. Hasil pemetaan antara ancaman, kebutuhan bukti digital, dan audit event yang dihasilkan ditunjukkan pada Tabel \ref{tab:threat-event-mapping}.

\begin{xltabular}
    {\textwidth}
    {|>{\centering\arraybackslash}p{1.8cm}
     |>{\raggedright\arraybackslash}p{5.1cm}
     |X|}

\caption{Pemetaan Ancaman ke Audit Event}
\label{tab:threat-event-mapping}\\

\hline
\textbf{Threat} &
\textbf{Bukti yang Dibutuhkan} &
\textbf{Audit Event} \\
\hline
\endfirsthead

\hline
\textbf{Threat} &
\textbf{Bukti yang Dibutuhkan} &
\textbf{Audit Event} \\
\hline
\endhead

\hline
\endfoot

\hline
\endlastfoot

T1, T2, T3 &
Identitas pengguna, metode autentikasi, hasil autentikasi, waktu, perangkat, serta frekuensi percobaan autentikasi. &
Authentication \\
\hline

T4, T5 &
Payload QR Access Delegation, hasil validasi, identitas pembuat dan pemindai QR. &
QR Access Delegation Validation \\
\hline

T6 &
Identitas pasien, penerima akses, \textit{capability} yang dibuat, waktu, serta \textit{transaction digest}. &
Capability Creation \\
\hline

T7, T13 &
Identitas pengguna, jenis operasi (\textit{read}/\textit{write}), objek rekam medis, \textit{capability} yang digunakan, serta hasil validasi otorisasi. &
Medical Record Access \\
\hline

T8 &
Identitas Hospital Administrator, penerima \textit{activation key}, dan waktu penerbitan. &
Activation Key Issuance \\
\hline

T9 &
Identitas Ministry Administrator, identitas fasyankes yang diregistrasi, waktu, serta \textit{transaction digest}. &
Healthcare Facility Registration \\
\hline

T10, T11, T12 &
Identitas pemohon, layanan PRE yang diminta, endpoint tujuan, waktu, dan hasil pemrosesan permintaan. &
PRE Service Request \\
\hline

T14 &
Identitas pemohon, aktivitas \textit{re-encryption}, objek yang diproses, dan hasil operasi. &
Re-encryption Operation \\
\hline

T15 &
\textit{Transaction digest}, identitas penandatangan, payload transaksi, dan hasil pengiriman transaksi. &
Blockchain Transaction Submission \\
\hline

T16, T18 &
Identitas pemohon sponsorship, \textit{transaction digest}, waktu, serta frekuensi permintaan sponsorship. &
Gas Sponsorship Request \\
\hline

T17 &
Identitas pemohon, keputusan sponsorship, dan \textit{transaction digest}. &
Gas Sponsorship Decision \\
\hline

T19 &
Jenis kunci yang dimodifikasi, jenis operasi, dan identitas pelaku. &
Client Key Store Operation \\
\hline

T20 &
Jenis kunci yang diakses, identitas pelaku, dan waktu akses. &
Client Key Store Access \\
\hline

T21, T23 &
Objek sensitif (\textit{nonce}/\texttt{kfrag}), jenis operasi, identitas proses, dan waktu operasi. &
Sensitive Cache Object Operation \\
\hline

T22 &
Objek sensitif yang diakses, identitas proses, dan waktu akses. &
Sensitive Cache Object Access \\
\hline

T24 &
Content Identifier (CID), identitas pemohon, jenis operasi, dan waktu akses. &
IPFS Object Access \\
\hline

T25 &
CID, nilai hash, hasil verifikasi integritas, dan waktu verifikasi. &
IPFS Object Integrity Verification \\
\hline

T26 &
Objek \textit{on-chain} yang diakses, identitas pemohon, dan waktu akses. &
On-chain Metadata Access \\
\hline

T27, T28 &
Capability ID, identitas pemohon, hasil validasi otorisasi, serta alasan penolakan apabila validasi gagal. &
Capability Validation \\
\hline

\end{xltabular}

Berdasarkan hasil pemetaan tersebut, diperoleh sekumpulan audit event yang mewakili aktivitas-aktivitas penting pada sistem DecMed. Satu audit event dapat digunakan sebagai bukti terhadap lebih dari satu ancaman apabila ancaman tersebut memiliki kebutuhan bukti digital yang serupa. Daftar audit event hasil identifikasi tersebut ditunjukkan pada Tabel \ref{tab:audit-event-list}.

\begin{xltabular}
    {\textwidth}
    {|>{\centering\arraybackslash}p{1cm}
     |>{\raggedright\arraybackslash}p{4cm}
     |X|}

\caption{Daftar Audit Event Hasil Identifikasi}
\label{tab:audit-event-list}\\

\hline
\textbf{ID} &
\textbf{Audit Event} &
\textbf{Deskripsi Aktivitas} \\
\hline
\endfirsthead

\hline
\textbf{ID} &
\textbf{Audit Event} &
\textbf{Deskripsi Aktivitas} \\
\hline
\endhead

\hline
\endfoot

\hline
\endlastfoot

EV1 & Authentication & Pencatatan setiap percobaan autentikasi pengguna ke Client, baik berhasil maupun gagal. \\
\hline

EV2 & QR Access Delegation Validation & Pencatatan proses validasi payload QR Access Delegation sebelum akses diberikan kepada tenaga fasyankes. \\
\hline

EV3 & Capability Creation & Pencatatan aktivitas pembuatan \textit{capability} akses oleh pasien kepada tenaga fasyankes. \\
\hline

EV4 & Medical Record Access & Pencatatan aktivitas pembacaan maupun perubahan data rekam medis. \\
\hline

EV5 & Activation Key Issuance & Pencatatan aktivitas penerbitan \textit{activation key} kepada personel fasyankes. \\
\hline

EV6 & Healthcare Facility Registration & Pencatatan aktivitas registrasi fasyankes beserta administratornya. \\
\hline

EV7 & PRE Service Request & Pencatatan setiap permintaan layanan yang diterima PRE Server. \\
\hline

EV8 & Re-encryption Operation & Pencatatan aktivitas re-enkripsi data oleh PRE Server. \\
\hline

EV9 & Blockchain Transaction Submission & Pencatatan aktivitas pembuatan, penandatanganan, dan pengiriman transaksi ke jaringan IOTA. \\
\hline

EV10 & Gas Sponsorship Request & Pencatatan aktivitas permintaan sponsorship transaksi kepada Gas Station. \\
\hline

EV11 & Gas Sponsorship Decision & Pencatatan keputusan Gas Station terhadap permintaan sponsorship transaksi. \\
\hline

EV12 & Client Key Store Operation & Pencatatan aktivitas pembuatan, perubahan, maupun penghapusan material kriptografi pada Client Key Store. \\
\hline

EV13 & Client Key Store Access & Pencatatan aktivitas akses terhadap material kriptografi pada Client Key Store. \\
\hline

EV14 & Sensitive Cache Object Operation & Pencatatan aktivitas penyimpanan maupun perubahan objek sensitif pada Redis. \\
\hline

EV15 & Sensitive Cache Object Access & Pencatatan aktivitas pembacaan objek sensitif pada Redis. \\
\hline

EV16 & IPFS Object Access & Pencatatan aktivitas pengunggahan maupun pengunduhan objek pada IPFS. \\
\hline

EV17 & IPFS Object Integrity Verification & Pencatatan proses verifikasi integritas objek yang tersimpan pada IPFS. \\
\hline

EV18 & On-chain Metadata Access & Pencatatan aktivitas pembacaan metadata maupun objek pada ledger IOTA. \\
\hline

EV19 & Capability Validation & Pencatatan proses validasi \textit{capability} sebelum permintaan diproses lebih lanjut. \\
\hline

\end{xltabular}

\subsubsection{Identifikasi Atribut Audit}

Setelah diperoleh daftar \textit{audit event} pada Tabel~\ref{tab:audit-event-list}, tahap selanjutnya adalah mengidentifikasi atribut yang perlu dicatat setiap kali \textit{audit event} tersebut terjadi. Kumpulan atribut tersebut membentuk sebuah \textit{audit record} yang menjadi bukti digital atas aktivitas sistem. Penentuan atribut audit dilakukan dengan mengacu pada kontrol AU-3 \textit{Content of Audit Records} pada NIST SP 800-53 Rev. 5 \cite{nist80053r5}, yang menyatakan bahwa setiap \textit{audit record} harus memuat informasi mengenai (a) jenis \textit{event} yang terjadi, (b) waktu kejadian, (c) lokasi atau sumber kejadian, (d) sumber \textit{event}, (e) hasil (\textit{outcome}) dari \textit{event}, serta (f) identitas individu, subjek, maupun objek yang terlibat.

Keenam elemen tersebut kemudian diimplementasikan dalam bentuk atribut audit umum yang dicatat pada seluruh \textit{audit event} sebagaimana ditunjukkan pada Tabel~\ref{tab:skema-umum-audit-record}. Selain atribut umum tersebut, setiap \textit{audit event} juga memiliki atribut tambahan yang disesuaikan dengan karakteristik aktivitas sistem dan kebutuhan bukti digital terhadap ancaman yang dipetakan sebelumnya. Atribut tambahan tersebut disajikan pada Tabel~\ref{tab:detail-audit-record}.

\begin{table}[h]
\centering
\caption{Skema Umum Atribut Audit}
\label{tab:skema-umum-audit-record}
\begin{tabular}{|>{\raggedright\arraybackslash}p{3.8cm}|>{\raggedright\arraybackslash}p{7.2cm}|c|}
\hline
\textbf{Atribut} & \textbf{Deskripsi} & \textbf{Elemen AU-3} \\
\hline
\texttt{record\_id} &
Pengenal unik setiap \textit{audit record} (UUID) yang tidak dapat digunakan ulang. &
-- \\
\hline

\texttt{event\_type} &
Jenis \textit{audit event} yang terjadi, mengacu pada daftar \textit{audit event} pada Tabel~\ref{tab:audit-event-list}. &
(a) \\
\hline

\texttt{timestamp} &
Waktu terjadinya \textit{audit event} dengan presisi milidetik dalam format UTC. &
(b) \\
\hline

\texttt{source\_component} &
Komponen atau proses sistem yang menghasilkan \textit{audit event}. &
(c), (d) \\
\hline

\texttt{actor\_id} &
Identitas aktor yang memicu \textit{audit event}, seperti \textit{IOTA address}, \textit{role}, atau \textit{session ID}. &
(f) \\
\hline

\texttt{target\_object} &
Identitas objek atau sumber daya yang menjadi target aktivitas, misalnya ID rekam medis, ID \textit{capability}, atau CID IPFS. &
(f) \\
\hline

\texttt{outcome} &
Hasil pelaksanaan \textit{audit event}, seperti \textit{success}, \textit{failure}, atau \textit{denied}, beserta kode alasan apabila tersedia. &
(e) \\
\hline

\texttt{prev\_record\_hash} &
\textit{Hash} dari \textit{audit record} sebelumnya untuk membentuk mekanisme \textit{hash chaining} yang mendukung sifat \textit{append-only} dan \textit{immutable}. &
-- \\
\hline

\end{tabular}
\end{table}

Field \texttt{record\_id} dan \texttt{prev\_record\_hash} tidak secara eksplisit disyaratkan oleh kontrol AU-3. Kedua atribut tersebut ditambahkan untuk memenuhi kebutuhan rancangan \textit{audit trail} pada penelitian ini, khususnya dalam menjamin identitas unik setiap \textit{audit record} serta menjaga integritas rangkaian \textit{audit record} melalui mekanisme \textit{hash chaining}.

\begin{xltabular}
{\textwidth}
{|>{\centering\arraybackslash}p{1.2cm}|X|>{\raggedright\arraybackslash}p{5.8cm}|}

\caption{Atribut Audit Tambahan untuk Setiap Audit Event}
\label{tab:detail-audit-record}\\

\hline
\textbf{ID Event} &
\textbf{Atribut Tambahan} &
\textbf{Alasan} \\
\hline
\endfirsthead

\hline
\textbf{ID Event} &
\textbf{Atribut Tambahan} &
\textbf{Alasan} \\
\hline
\endhead

\hline
\endfoot

\hline
\endlastfoot

EV1 &
\texttt{auth\_method}, \texttt{authentication\_result}, \texttt{failed\_attempt\_count}, \texttt{device\_fingerprint}. &
Mendukung identifikasi metode autentikasi, pola kegagalan autentikasi, serta anomali perangkat yang digunakan. \\
\hline

EV2 &
\texttt{qr\_payload\_id}, \texttt{recipient\_identity}, \texttt{signature\_valid}. &
Membuktikan keaslian payload QR serta identitas pihak yang menerima delegasi akses. \\
\hline

EV3 &
\texttt{capability\_id}, \texttt{access\_scope}, \texttt{expiry\_duration}, \texttt{transaction\_digest}. &
Mendokumentasikan capability yang diberikan beserta ruang lingkup akses dan transaksi yang merekamnya. \\
\hline

EV4 &
\texttt{access\_type}, \texttt{medical\_record\_id}, \texttt{capability\_id}, \texttt{authorization\_token\_id}. &
Mengidentifikasi jenis akses terhadap rekam medis beserta capability dan token otorisasi yang digunakan. \\
\hline

EV5 &
\texttt{activation\_key\_id}, \texttt{issuer\_id}, \texttt{recipient\_id}, \texttt{expiration\_time}. &
Mencatat informasi penerbitan \textit{activation key} beserta pihak penerbit dan penerimanya. \\
\hline

EV6 &
\texttt{facility\_id}, \texttt{facility\_name}, \texttt{administrator\_id}, \texttt{transaction\_digest}. &
Mendokumentasikan registrasi fasilitas pelayanan kesehatan beserta transaksi yang merekamnya. \\
\hline

EV7 &
\texttt{endpoint\_called}, \texttt{request\_id}, \texttt{caller\_component}, \texttt{channel\_encryption}. &
Mengidentifikasi layanan PRE yang dipanggil beserta komponen pemanggil dan status pengamanan kanal komunikasi. \\
\hline

EV8 &
\texttt{reencryption\_operation\_id}, \texttt{capability\_id}, \texttt{target\_ciphertext}, \texttt{kfrag\_identifier}. &
Mendokumentasikan proses re-enkripsi yang dilakukan terhadap data terenkripsi. \\
\hline

EV9 &
\texttt{transaction\_digest}, \texttt{payload\_hash}, \texttt{network\_confirmation\_status}. &
Menyediakan bukti pengiriman transaksi beserta status konfirmasi pada jaringan IOTA. \\
\hline

EV10 &
\texttt{requested\_gas\_budget}, \texttt{requester\_id}, \texttt{transaction\_digest}. &
Mendokumentasikan permintaan sponsorship beserta kebutuhan gas transaksi. \\
\hline

EV11 &
\texttt{approval\_status}, \texttt{decision\_reason}, \texttt{approved\_by}. &
Mencatat keputusan sponsorship beserta alasan keputusan dan pihak yang memberikan persetujuan. \\
\hline

EV12 &
\texttt{key\_operation}, \texttt{key\_id}, \texttt{key\_type}. &
Mencatat aktivitas pembuatan, perubahan, maupun penghapusan material kriptografi. \\
\hline

EV13 &
\texttt{key\_id}, \texttt{key\_type}, \texttt{access\_purpose}. &
Mencatat aktivitas akses terhadap material kriptografi beserta tujuan penggunaannya. \\
\hline

EV14 &
\texttt{redis\_key\_type}, \texttt{operation\_type}, \texttt{ttl\_remaining}. &
Mencatat aktivitas penyimpanan atau perubahan objek sensitif pada Redis. \\
\hline

EV15 &
\texttt{redis\_key\_type}, \texttt{request\_origin}, \texttt{ttl\_remaining}. &
Mencatat aktivitas pembacaan objek sensitif beserta asal permintaan. \\
\hline

EV16 &
\texttt{cid}, \texttt{operation\_type}, \texttt{data\_size}. &
Mencatat aktivitas unggah maupun unduh objek pada penyimpanan IPFS. \\
\hline

EV17 &
\texttt{cid}, \texttt{expected\_hash}, \texttt{verification\_result}. &
Mendokumentasikan proses verifikasi integritas objek pada IPFS. \\
\hline

EV18 &
\texttt{metadata\_type}, \texttt{object\_id}, \texttt{query\_requester\_id}. &
Mencatat aktivitas pembacaan metadata atau objek pada ledger IOTA. \\
\hline

EV19 &
\texttt{capability\_checked}, \texttt{required\_scope}, \texttt{actual\_scope}, \texttt{validation\_result}. &
Mendokumentasikan proses validasi capability beserta hasil evaluasi hak akses. \\
\hline

\end{xltabular}

\subsubsection{Analisis Gap Audit Event}

Setelah diperoleh daftar \textit{audit event} yang dibutuhkan, tahap selanjutnya adalah mengevaluasi apakah sistem DecMed telah menyediakan mekanisme pencatatan terhadap \textit{audit event} tersebut. Evaluasi dilakukan dengan membandingkan daftar \textit{audit event} hasil identifikasi dengan mekanisme pencatatan yang telah tersedia pada DecMed, yaitu \textit{PatientAccessLog} yang disimpan sebagai objek \textit{on-chain} serta riwayat transaksi (\textit{transaction history}) pada jaringan IOTA. Hasil analisis tersebut disajikan pada Tabel~\ref{tab:audit-event-gap}.

\begin{xltabular}
{\textwidth}
{|>{\centering\arraybackslash}p{1cm}
 |>{\raggedright\arraybackslash}p{2cm}
 |>{\centering\arraybackslash}p{2cm}
 |X|}

\caption{Analisis Gap Audit Event pada DecMed}
\label{tab:audit-event-gap}\\

\hline
\textbf{ID} &
\textbf{Audit Event} &
\textbf{Kondisi Saat Ini} &
\textbf{Analisis Gap} \\
\hline
\endfirsthead

\hline
\textbf{ID} &
\textbf{Audit Event} &
\textbf{Kondisi Saat Ini} &
\textbf{Analisis Gap} \\
\hline
\endhead

\hline
\endfoot

\hline
\endlastfoot

EV1 &
Authentica\-tion &
Tidak tersedia &
Tidak terdapat mekanisme pencatatan aktivitas autentikasi pengguna sehingga percobaan autentikasi berhasil maupun gagal tidak dapat dijadikan bukti audit. \\
\hline

EV2 &
QR Access Delegation Validation &
Tidak tersedia &
Proses validasi payload QR tidak menghasilkan audit record sehingga tidak tersedia bukti terhadap penyalahgunaan maupun manipulasi QR Access Delegation. \\
\hline

EV3 &
Capability Creation &
Sebagian &
Aktivitas tercatat pada \textit{PatientAccessLog} dan transaksi \textit{on-chain}, namun informasi yang dicatat masih terbatas dan belum dirancang sebagai audit trail yang terintegrasi. \\
\hline

EV4 &
Medical Record Access &
Tidak tersedia &
Aktivitas pembacaan maupun perubahan data rekam medis belum dicatat sebagai audit record sehingga sulit dilakukan penelusuran akses terhadap data pasien. \\
\hline

EV5 &
Activation Key Issuance &
Tidak tersedia &
Tidak tersedia mekanisme audit terhadap proses penerbitan \textit{activation key}. \\
\hline

EV6 &
Healthcare Facility Registration &
Sebagian &
Aktivitas registrasi direkam sebagai transaksi \textit{on-chain}, namun belum tersedia audit record yang mendokumentasikan informasi kontekstual aktivitas registrasi. \\
\hline

EV7 &
PRE Service Request &
Tidak tersedia &
Aktivitas permintaan layanan pada PRE Server belum dicatat sebagai audit record. \\
\hline

EV8 &
Re-encryption Operation &
Tidak tersedia &
Proses \textit{re-encryption} tidak menghasilkan bukti audit yang dapat digunakan dalam investigasi. \\
\hline

EV9 &
Blockchain Transaction Submission &
Sebagian &
Riwayat transaksi tersedia pada jaringan IOTA, namun hanya merekam hasil transaksi tanpa informasi operasional yang lengkap mengenai proses pengiriman transaksi. \\
\hline

EV10 &
Gas Sponsorship Request &
Tidak tersedia &
Tidak tersedia mekanisme audit terhadap permintaan sponsorship transaksi. \\
\hline

EV11 &
Gas Sponsorship Decision &
Tidak tersedia &
Keputusan pemberian maupun penolakan sponsorship belum dicatat sebagai audit record. \\
\hline

EV12 &
Client Key Store Operation &
Tidak tersedia &
Aktivitas pengelolaan material kriptografi pada sisi klien belum memiliki mekanisme audit. \\
\hline

EV13 &
Client Key Store Access &
Tidak tersedia &
Akses terhadap material kriptografi tidak menghasilkan bukti audit. \\
\hline

EV14 &
Sensitive Cache Object Operation &
Tidak tersedia &
Aktivitas penyimpanan objek sensitif pada Redis belum dicatat. \\
\hline

EV15 &
Sensitive Cache Object Access &
Tidak tersedia &
Aktivitas pembacaan objek sensitif pada Redis belum dicatat. \\
\hline

EV16 &
IPFS Object Access &
Tidak tersedia &
Aktivitas unggah maupun unduh objek pada IPFS belum memiliki mekanisme audit. \\
\hline

EV17 &
IPFS Object Integrity Verification &
Tidak tersedia &
Tidak tersedia pencatatan terhadap proses verifikasi integritas objek pada IPFS. \\
\hline

EV18 &
On-chain Metadata Access &
Tidak tersedia &
Aktivitas pembacaan metadata maupun objek \textit{on-chain} tidak dicatat sebagai audit record. \\
\hline

EV19 &
Capability Validation &
Tidak tersedia &
Proses validasi \textit{capability} sebelum otorisasi akses belum menghasilkan audit record. \\
\hline

\end{xltabular}

Berdasarkan hasil analisis pada Tabel~\ref{tab:audit-event-gap}, dapat diketahui bahwa DecMed telah menyediakan mekanisme pencatatan terhadap sebagian aktivitas melalui \textit{PatientAccessLog} dan riwayat transaksi pada jaringan IOTA. Namun, mekanisme tersebut hanya mencakup sebagian kecil \textit{audit event} yang dibutuhkan serta belum dirancang sebagai sebuah sistem \textit{audit trail}. Sebagian besar aktivitas penting, khususnya yang terjadi pada komponen \textit{Client}, PRE Server, Redis, maupun IPFS, belum menghasilkan bukti audit yang dapat digunakan dalam proses audit maupun investigasi insiden keamanan. Selain itu, informasi yang tersedia pada mekanisme pencatatan yang telah ada masih bersifat terbatas dan tersebar pada beberapa komponen sehingga belum mampu menyediakan jejak audit yang lengkap. Oleh karena itu, diperlukan mekanisme audit trail yang mampu mencatat seluruh \textit{audit event} yang telah diidentifikasi secara terintegrasi.